% ECONOMICS_PROBLEM: {"id": "L13-382", "title": "Common ownership mechanisms", "jel_code": "L13", "source_id": "OP-0397"}
% !TeX program = xelatex
\documentclass[11pt,a4paper]{article}
\usepackage[margin=23mm]{geometry}
\usepackage{fontspec}
\setmainfont{Latin Modern Roman}
\usepackage{amsmath,amssymb,mathtools}
\usepackage{xcolor,enumitem,booktabs,longtable,array,xurl,hyperref}
\definecolor{muted}{HTML}{53636C}
\hypersetup{colorlinks=true,urlcolor=blue,linkcolor=blue}
\setlength{\parindent}{0pt}
\setlength{\parskip}{5pt}
\newcommand{\R}{\mathbb R}
\newcommand{\E}{\mathbb E}
\newcommand{\N}{\mathbb N}
\newcommand{\ind}{\mathbf 1}
\newcommand{\Var}{\operatorname{Var}}
\newcommand{\Cov}{\operatorname{Cov}}
\newcommand{\supp}{\operatorname{supp}}
\DeclareMathOperator*{\argmin}{arg\,min}
\DeclareMathOperator*{\argmax}{arg\,max}
\newcommand{\entrymeta}[3]{{\sffamily\small\color{muted}\textbf{\texttt{#1}}\quad|\quad #2\par #3\par}}
\newcommand{\Problemref}[1]{\texttt{#1}}
\newcommand{\JELref}[2]{\texttt{#2} (JEL \texttt{#1})}
\begin{document}
\section*{Common ownership mechanisms}
\textbf{Problem ID:} \texttt{L13-382}\quad \textbf{JEL:} \texttt{L13}\par
% BEGIN ECONOMICS STATEMENT
\label{problem:OP-0397}
{\sffamily\small\color{muted}Original subject group: Industrial organization and competition\par}
\label{definitions:problem:30007679}
\textbf{Audit classification:} Published research agenda. The common-ownership passage was verified. Ownership-induced governance is a causal mediator, not an automatically observed exogenous regressor.
\subsection*{Definitions and precise research target}
Consider publicly traded United States grocery manufacturers observed annually for ten years. Let \(O_t\) be the matrix of institutional investors' voting shares in these firms, and let \(G_t\) record investor engagements, compensation provisions and board decisions. Define \(p_j(O,G,x)\) as firm \(j\)'s equilibrium price under ownership \(O\), governance arrangement \(G\), and market fundamentals \(x\), including demand and production costs. For a prespecified ownership change from \(O^0\) to \(O^1\), target the total price effect \(E[p_j(O^1,G(O^1),x)-p_j(O^0,G(O^0),x)]\), where \(G(O)\) is governance induced by ownership and expectation averages over products and markets. Also target the controlled effect of ownership with governance held at a fixed feasible \(G^0\). Identify a source of ownership variation independent of latent competitive prospects, and specify how investor portfolio changes can alter governance rather than assuming portfolio-profit maximization by managers. Mediated effects require additional assumptions about ownership-induced confounders of governance and pricing; report ranges when these assumptions are unsupported. Compare predicted mechanisms with engagement and incentive records. The task seeks causal magnitude and transmission channels in this population, separating common ownership from direct cross-holdings and common shocks. Observed ownership concentration alone does not identify softer product-market rivalry.

\textbf{Definitions, assumptions and mathematical qualifications.} Ownership entries \(O_{hj}\in[0,1]\) are voting shares, with \(\sum_hO_{hj}\leq1\); common ownership means an investor holds shares in two competing firms, distinct from firms owning each other. A structural governance law \(G(O,U_G)\) and price law \(p_j(O,G,x,U_p)\) use a common exogenous-shock distribution across interventions. Total and controlled effects are respectively \(E[p_j(O^1,G(O^1,U_G),x,U_p)-p_j(O^0,G(O^0,U_G),x,U_p)]\) and the same contrast with \(G=G^0\). The intervention must make \(G^0\) feasible under both ownership matrices. Specify weights over products and markets, an equilibrium-selection kernel and the confounders controlled by the ownership design.
\subsection*{Variants and assumption changes}
\begin{enumerate}
\item \textbf{Editorial, controlled intervention.} Randomize or instrument a feasible governance provision while holding ownership fixed, and estimate its price effect under explicit exclusion and relevance assumptions.
\item \textbf{Editorial, mediation.} Target a natural indirect effect only after stating cross-regime mediator coupling and no ownership-induced mediator--outcome confounding; otherwise bound total and controlled effects without assigning a unique mediated component.
\end{enumerate}
\subsection*{References and source audit}
\textbf{Checked source.} 24 August 2024 author draft, Section 5, printed pp. 33--34, common-ownership discussion. Full primary text retrieved. \url{https://web.stanford.edu/~ayurukog/trendsincompetition.pdf}.
\begin{enumerate}\raggedright
\item\hypertarget{definitions:source:30007679:1}{} Carl Shapiro; Ali Yurukoglu. 2024. \emph{Trends in Competition in the United States: What Does the Evidence Show?}. 24 August 2024 author draft, Section 5, printed pp. 33--34, common-ownership discussion.
\url{https://web.stanford.edu/~ayurukog/trendsincompetition.pdf}.
DOI: \url{https://doi.org/10.3386/w32762}.
\end{enumerate}

\subsection*{Common conventions: Source mathematical conventions}

These conventions apply to each entry unless explicitly replaced.

\textbf{Models and identification.} A primitive model \(m\) specifies all probability laws, feasible choices, information, equilibrium and measurement rules used by its target. The admissible class \(\mathcal M\) is fixed independently of outcomes. If \(P_O(m)\) is its observable probability law and \(\tau(m)\) the target, its identified set at observable law \(P\) is
\[
 \mathcal I_\tau(P)=\{\tau(m):m\in\mathcal M,\ P_O(m)=P\}.
\]
Point identification means that this set is a singleton. An empty set signals incompatibility of data and assumptions. Coordinate bounds need not describe the joint set. Bounds are sharp only if every retained target is attainable or arbitrarily approachable by compatible models. Finite bounds require explicit restrictions on unobserved quantities; unrestricted confounding, hidden wealth, extrapolation or arbitrary equilibrium selection can yield uninformative sets.

\textbf{Causal interventions.} A potential outcome \(Y_i(p)\) is the outcome under a complete policy regime \(p\), including timing, financing, information and market spillovers. The target population and units are fixed before treatment. With interference, use the complete assignment vector or a justified exposure mapping; individual no-interference notation alone cannot identify an economy-wide intervention. A difference of mean potential outcomes does not require the cross-world joint distribution. Conditioning on post-treatment migration, survival, enrollment or employment changes the estimand and needs separate justification.

\textbf{Statistical statements.} An estimator, confidence set, forecast or test specifies a sampling law, model class and loss/coverage criterion. Uniform \((1-\alpha)\)-coverage means \(\inf_{m\in\mathcal M}\Pr_m\{\tau(m)\in C_n\}\geq1-\alpha\), or an explicitly asymptotic version. The whole parameter space trivially has coverage, so informativeness requires a stated width, risk or power objective. A forecasting improvement is relative to a named benchmark, information set and evaluation population.

\textbf{Equilibrium and optimization.} An equilibrium comprises feasible plans, optimality, beliefs where needed, budget constraints and all market/resource clearing. The primitives or a stated benchmark define each correspondence \(\mathcal E(p;m)\). Nonemptiness is assumed or proved, never inferred just from compact policy sets. A selected-equilibrium target uses a declared selection rule; a robust target quantifies over all permitted equilibria. Use suprema or infima unless attainment is established. An \(\varepsilon\)-optimal policy is feasible and within \(\varepsilon\) of the finite optimum value. Report an empty feasible set separately.

\textbf{Welfare and accounting.} Normative utility, interpersonal normalization, social weights, numeraire, discounting and the use of public receipts are primitives. Behavioral utility need not coincide with the welfare criterion. Household welfare includes ownership income and rents once; adding profits again to compensating variation generally double-counts them. A monetary equivalent is defined by an explicit transfer and price convention, with monotonicity and range conditions. Average effects, mechanism contributions, transfers and welfare losses are different targets. Infinite sums require convergence and dynamic feasible plans require terminal or no-Ponzi conditions.

\textbf{Source versions.} A working-paper date, revision date and journal publication year are distinguished. A DOI for a journal version is not automatically the DOI of an earlier manuscript. Locators refer to the stated version and printed pages where specified. A metadata-only check does not verify a claim about a full-text conclusion. Every entry's source subsection narrows the provenance claim accordingly.
% END ECONOMICS STATEMENT
\end{document}
